\subsection{Selection of papers for the detailed review in this report}
Two samples of fourteen papers each are individually reviewed in
Section~\ref{sec:results-detail}: a classical-computing sample selected from the CS-side
review queue described above, and a quantum-computing sample selected from the
quantum review queue, chosen to span a comparable range of outcomes ---
directly-linked data, named-but-unlinked data, request-conditional access, and
unresolvable or not-yet-live links --- so that the two samples are a fair basis for
qualitative comparison despite their small size. Each paper, in both samples, was
processed independently as follows.

\subsection{Extraction, link verification, and classification}
\label{sec:extraction-classification}

Once a paper is selected from a domain's corpus (Section~\ref{sec:candidate-list-gen}),
it passes through extraction and live link verification, classification, and
review-CSV construction before it appears in the aggregate results
(Section~\ref{sec:results-aggregate}). These steps are implemented as two scripts plus
one non-scripted classification step, described in turn below.

\paragraph{Extraction and live link verification (scripted).} A dedicated
script reuses the plain text already cached by the link-discovery step
(Section~\ref{sec:candidate-list-gen}) wherever available, avoiding a second download
of the same PDF; a paper not previously scanned is downloaded and extracted fresh using
the same procedure. Every link previously extracted for that paper is then
independently re-checked over HTTP (a HEAD request, falling back to a small ranged GET
for servers that reject HEAD), and classified as resolving, returning an error, or
timing out. Every paper in a domain's corpus goes through this step, so a paper that
uses external data without ever using an availability keyword or including a link is
still fully re-checked here. The
result of this step is one pending record per paper, containing the full extracted
text, the live link-verification results, and the classification fields left unset.

\paragraph{Batching for full-corpus scale (scripted preprocessing, not classification
itself).} At full-corpus scale, a paper's extracted text is too large to read in bulk
without either truncating it or exceeding practical per-request limits (a single
paper's text ranges from roughly 25{,}000 to over 500{,}000 characters). Two scripted
steps address this without changing what is actually classified. First, the pending
records for a batch of papers are condensed into a digest: papers short enough to read
whole are left as-is; longer papers are reduced to their opening excerpt (title,
abstract, and introduction, where a paper's external data sources are typically first
named), their closing excerpt (references and, commonly, a formal data-availability
statement), and a window of text around every occurrence of a data/code-availability
keyword found elsewhere in the paper, so a relevant sentence in the middle of a long
paper is not missed by only keeping the ends. Classification is performed from this
digest plus the unmodified link-verification results, on the same basis it would be
performed from the full text --- the digest changes how much text is read, not what is
asked of it or how it is judged. Second, once a batch is classified, a separate script
writes the classification fields onto a copy of the paper's original, complete record
(full extracted text and link-verification results, unmodified), so the digest is
never what persists in the cache; it is a reading aid discarded after use.

\paragraph{Classification (agent-driven, not a scripted API call).} The two-axis,
six-case framework (Section~\ref{sec:framework}) is defined in exactly one place --- a
single framework document --- and every classification, whether performed by a person
or by an automated agent, is required to read that same document rather than work from
a separately maintained copy, so that the framework's definitions cannot drift out of
sync between where they are documented and how they are actually applied. Concretely,
classification in this project is performed by a Claude Code agent: for each pending
record, the agent reads the framework document together with that paper's extracted
text and link-verification results, and assigns an Axis A case, an Axis B case, a
short evidence quote for each, a confidence level, and a flag indicating whether the
assignment should be treated as provisional pending future human confirmation. This
step is deliberately not implemented as a scripted, fire-and-forget call to a hosted
language-model API: the extraction script stops at extraction and verification only,
and classification instead runs as a supervised agent process --- one that is paused,
corrected, and resumed batch by batch, rather than an unattended pipeline with no point
of judgment in the loop. As of this writing, every classification produced this way has
been performed by the agent alone.

\paragraph{Distinguishing the paper's own data/code link from an unrelated citation.}
A paper can contain many links that have nothing to do with its own data or code ---
a citation to a library it used, a link to a related paper, a reference to a
co-author's unrelated prior work. Link discovery (Section~\ref{sec:candidate-list-gen})
does not filter for this: it collects every candidate link, DOI, and URL in the paper
regardless of what it points to. Deciding which of those links, if any, is actually
the paper's own data/code artifact is not a separate rule-based step; it is made by
the classifying agent as part of the same read that assigns each case, using the same
surrounding text a human reviewer would rely on: a link's ownership is judged from the
sentence or paragraph it appears in (``our code is available at...'' versus ``we used
the Qiskit library [12]''), from whether it sits under an explicit data/code-
availability statement, and from whether the paper's own text claims the linked
resource as its own contribution rather than citing it as a tool or related work. No
separate attribution rule exists in the pipeline for this judgment; it is treated as
ordinary reading comprehension, the same as a human reviewer would apply, rather than
a distinct classification rule.

\paragraph{Review CSV construction (scripted).} A second script joins the
classification output back against the domain's paper-list metadata (title, authors,
publication date) and writes one CSV per domain, with human-readable case labels,
the evidence quotes, the link-verification outcomes, and a blank verification column.

\paragraph{Sequential batch processing.} Batches are classified sequentially rather
than concurrently: this trades wall-clock time for the property that a failure
partway through affects only the batch in progress rather than all of them at once,
and work already merged from completed batches is retained regardless of what happens
to a later one.

\paragraph{Human verification.} No agent-produced classification is treated as
final on its own. Every row is written with its verification column blank; a human
reviewer subsequently confirms or overrides each case assignment, following the same
convention already in use elsewhere in this project's review queues (verification
status recorded separately from the label itself, so the two are never conflated).
This mirrors the project's broader design principle, applied consistently from the
candidate-generation stage onward: an ambiguous or low-confidence signal is flagged for
review rather than silently resolved one way or the other.
